% Generated draft: requires mathematical and source review.
\documentclass[11pt]{article}
\usepackage[margin=25mm]{geometry}
\usepackage{fontspec}
\setmainfont{Latin Modern Roman}
\usepackage{amsmath,amssymb,mathtools,unicode-math}
\setmathfont{Latin Modern Math}
\usepackage{enumitem,needspace,xurl,hyperref}
\setlength{\parindent}{0pt}
\setlength{\parskip}{5pt}
\setlength{\emergencystretch}{3em}
\newcommand{\N}{\mathbb{N}}
\newcommand{\Z}{\mathbb{Z}}
\newcommand{\Q}{\mathbb{Q}}
\newcommand{\R}{\mathbb{R}}
\newcommand{\C}{\mathbb{C}}
\newcommand{\F}{\mathbb{F}}
\newcommand{\A}{\mathbb{A}}
\newcommand{\PP}{\mathbb{P}}
\newcommand{\E}{\mathbb{E}}
\newcommand{\eps}{\varepsilon}
\newcommand{\dd}{\,\mathrm{d}}
\begin{document}
\section*{UnsolvedMath ID: 20003138}
\textbf{Source problem code:} catalogue entry
\subsection{Definitions and mathematical statement}
The source asks whether a noncollapsing Gromov--Hausdorff convergence $M_i\to X$ carries a ``tangent bundle'' structure on $X$ sensitive to the diffeomorphism type of the approximating manifolds. It specifies neither curvature/compactness assumptions nor what sensitivity or bundle means. The workshop's nonnegative-sectional-curvature setting is relevant background, but noncollapse alone does not provide the homeomorphisms needed below.

The aid restricts to dimension $n\geq5$ and the standard Perelman-stability regime: for example, closed $n$-manifolds with a common lower sectional-curvature bound, bounded diameter, noncollapsing volume and compact $n$-dimensional Alexandrov limit. Stability supplies homeomorphisms between $X$ and the sufficiently large $M_i$, so $X$ is a topological manifold in this regime. These are explicit additional hypotheses, not a theorem for every noncollapsed metric limit.

Let $\tau_{\mathrm{TOP}}:X\to BTOP$ classify the stable topological tangent microbundle, and let $j:BO\to BTOP$ forget vector-bundle linear structure. A stable $O$-reduction with its comparison data is a pair
\[
 \lambda:X\to BO,\qquad H:j\circ\lambda\simeq\tau_{\mathrm{TOP}}.
\]
The marked datum is the homotopy class of $(\lambda,H)$ in the homotopy fiber over $\tau_{\mathrm{TOP}}$ of
\[
 \operatorname{Map}(X,BO)\longrightarrow\operatorname{Map}(X,BTOP).
\]
It retains the homotopy $H$; the stable real vector-bundle class $[\lambda]$ alone is less information.

For a stability marking $h_i:X\to M_i$, the aid reports that the smooth tangent bundle $h_i^*TM_i$, together with the induced identification of its underlying topological bundle with the tangent microbundle of $X$, gives such a pair $(\lambda_i,H_i)$. Changing the marking acts by $\operatorname{Homeo}(X)$, including its induced tangential comparison. The corresponding orbit is reported independent of the marking, and distinct orbits obstruct diffeomorphism of the approximants. Orientation-preserving markings and homeomorphisms must be used if oriented diffeomorphism types are to be distinguished.

The summary further states that sufficiently controlled-close stability markings give the same marked homotopy-lift class. The exact control, compactness and isotopy theorem behind this assertion are not supplied. The construction uses smooth approximation data and a stability marking; it is not an invariant determined by the metric space $X$ alone, nor an assertion that every arbitrary choice of homeomorphism yields the same marked class.

The aid emphasizes a hierarchy of insufficient substitutes: metric tangent cones, the topological tangent microbundle, and even the underlying stable vector-bundle class can forget smoothing information. Its example is the 28 oriented smoothings of $S^7$, with identical stably trivial tangent-bundle classes. The comparison homotopies can still carry different smoothing data. The oriented count is not asserted as 28 distinct unoriented types, and the record does not construct a common nonnegative-curvature limit realizing all these smoothings. This is a smoothing-theoretic partial repair in the stated stable range, not a general construction for singular limits or dimensions below five.
\subsection{Short English statement}
The reported structure retains a stable vector-bundle reduction together with its comparison to the topological tangent bundle. Forgetting that comparison can lose exotic smoothings. The construction depends on stability hypotheses and the approximating smooth manifolds, rather than on the limit metric alone.
\subsection{Sources}
\begin{itemize}
\item Local UnsolvedMath registry, permanent ID 20003138, source code catalogue entry; \texttt{aim-topology-notes.json}, notes[225].
\item AIM, \emph{Manifolds with nonnegative sectional curvature}, Problem 13: \url{https://aimath.org/WWN/nnsectcurvature/nnsectcurvature.pdf}. No proposer, remarks or original literature entry is supplied.
\item Perelman stability, stable smoothing/$O$-reduction theory, controlled isotopy and the 28 oriented $S^7$ smoothings are invoked by the supplied aid without full bibliographic entries or exact theorem statements.
\item Machine-generated research attempt 1: partial result, partially solved, rigor labelled ``proved,'' low novelty confidence. Not peer reviewed; independently verify before citation. Registry updated 14 August 2026.
\end{itemize}
\end{document}
